\chapter{DPF HMC Target Correctness and Structural-Model Suitability}
\label{ch:bf-dpf-hmc-target-suitability}

The preceding DPF chapters defined a ladder of increasingly ambitious
constructions:
\begin{enumerate}[label=(\roman*)]
  \item classical particle filtering;
  \item particle-flow transport;
  \item PF-PF proposal correction;
  \item differentiable resampling and OT relaxation;
  \item learned or amortized OT as a further approximation layer.
\end{enumerate}
This chapter asks the question that none of the earlier chapters could answer on
its own: if one uses one of these constructions inside Hamiltonian Monte Carlo,
what target is HMC actually sampling, and when is that target interpretation
acceptable for nonlinear DSGE and MacroFinance models?

The central lesson is that HMC correctness is not a property of a method label
such as ``differentiable particle filter.''  It is a property of a concrete
value-gradient construction.  If the scalar log target and the gradient supplied
to the integrator correspond to different mathematical objects, then even a
smooth and numerically stable chain is exploring the wrong Hamiltonian system.

\section{DPF target contract for HMC}
\label{sec:bf-dpf-hmc-contract}

Let $\ell(\theta)$ denote the scalar log target supplied to HMC and let
$g(\theta)$ denote the gradient used by the numerical integrator.  The minimum
contract is
\begin{equation}
\label{eq:bf-dpf-hmc-contract}
  g(\theta) = \nabla_\theta \ell(\theta)
\end{equation}
for the same mathematical object $\ell$.  This does not force $\ell$ to be an
exact nonlinear filtering likelihood.  It does force the user to say what
object $\ell$ is.

For the DPF ladder, the relevant target statuses are:
\begin{enumerate}[label=(\alph*)]
  \item exact target in a special case;
  \item unbiased particle-based likelihood estimator used inside a corrected
    MCMC interpretation;
  \item approximate or relaxed target;
  \item learned surrogate target.
\end{enumerate}
The first task of this chapter is to place each DPF rung into one of these
categories.

\section{Rung-by-rung target analysis}
\label{sec:bf-dpf-hmc-rung-analysis}

\subsection{EDH}

At the EDH rung, the scalar is a flow-based approximate likelihood object built
from the Gaussian-closure transport construction.  If the same object is
 differentiated, then the value and gradient are internally consistent, but the
object is generally not the exact nonlinear filtering likelihood outside the
linear-Gaussian special case.  Therefore EDH supports:
\begin{itemize}
  \item an exact-target interpretation only in the linear-Gaussian recovery
    benchmark;
  \item otherwise an approximate-target or value-side benchmark interpretation.
\end{itemize}
Smoothness of the flow ODE does not by itself upgrade EDH to a generic exact
HMC target.

\subsection{EDH/PF and LEDH/PF with proposal correction}

At the PF-PF rung, the scalar becomes a proposal-corrected particle quantity.
The correction restores the proposal-to-target density ratio through the
change-of-variables Jacobian and the importance weight.  If the same corrected
object is evaluated and differentiated, then this rung is the first serious
DPF-based HMC candidate in the BayesFilter ladder.

The remaining caution is not conceptual but quantitative.  Even when the target
interpretation is clearer than raw flow transport, the resulting HMC path may
still be fragile because of:
\begin{itemize}
  \item finite-particle variability;
  \item flow discretization error;
  \item Jacobian or log-determinant numerical error;
  \item compiled execution instability in repeated long-run evaluation.
\end{itemize}
Thus PF-PF is the first rung with a strong target story, but not automatically a
production-ready HMC path.

\subsection{Soft differentiable resampling}

At the soft-resampling rung, the scalar is built from a modified resampling map.
If HMC uses this scalar and its gradient consistently, the resulting chain may
still be mathematically coherent, but it should be interpreted as targeting the
soft-resampling surrogate posterior, not the exact categorical-resampling
posterior.  The gain is pathwise differentiability; the cost is target drift.

The key point is not that soft resampling is unusable.  The key point is that it
must be classified honestly.  Its mathematical status is surrogate-target HMC,
not exact-target HMC.

\subsection{OT and entropic OT resampling}

At the OT rung, the scalar is built from a transport-relaxed resampling map.
Relative to soft resampling, the geometric structure is cleaner, because the
relaxation is tied to an explicit transport problem.  Yet the target remains a
relaxed one: the entropic regularization and finite Sinkhorn solve change the
resampling law relative to exact categorical resampling.

Therefore, if the OT-relaxed scalar and its gradient are used consistently, the
appropriate interpretation is approximate-target or relaxed-target HMC.  This
may be a respectable and useful target, but it should not be mistaken for the
exact categorical-resampling target.

\subsection{Learned or amortized OT}

At the learned-OT rung, one introduces a further approximation to the OT
baseline.  The resulting HMC path is therefore a surrogate built on top of a
relaxed target.  If the learned scalar and its gradient are used consistently,
then the HMC path may again be mathematically coherent, but its target status is
farthest from the classical exact categorical-resampling interpretation.

The lesson is not that learned OT cannot be useful.  The lesson is that it must
be read as a second approximation layer, not merely as a cheap exact OT solver.

\section{Summary table}
\label{sec:bf-dpf-hmc-summary}

\begin{table}[ht]
\centering
\small
\begin{tabular}{p{0.16\textwidth} p{0.20\textwidth} p{0.19\textwidth} p{0.28\textwidth}}
\hline
Rung & Evaluated object & Differentiated object & HMC interpretation \\
\hline
EDH & flow-based approximate likelihood object & derivative of the same approximate flow object & exact only in the linear-Gaussian recovery benchmark; otherwise approximate-target benchmark \\
PF-PF & proposal-corrected particle estimator & derivative of the same corrected estimator & first serious DPF HMC candidate, still limited by Monte Carlo variability and numerical Jacobian issues \\
Soft resampling & relaxed differentiable resampling target & derivative of the same relaxed target & surrogate-target HMC \\
OT / EOT & transport-relaxed target & gradient through the same OT-relaxed target & approximate/relaxed-target HMC \\
Learned OT & learned approximation to OT-relaxed target & derivative of the same learned surrogate & surrogate-target HMC with an additional learned approximation layer \\
\hline
\end{tabular}
\caption{Rung-by-rung target interpretation for the DPF ladder.}
\label{tab:bf-dpf-hmc-rungs}
\end{table}

Table~\ref{tab:bf-dpf-hmc-rungs} is the main output of this chapter.  It makes
explicit why differentiability alone is insufficient: two methods may both yield
smooth gradients, yet those gradients may correspond to mathematically different
targets.

\section{Why nonlinear DSGE models sharpen the issue}
\label{sec:bf-dpf-hmc-dsge}

Nonlinear DSGE models make target drift especially important because the
state-space law already depends on structural timing, determinacy restrictions,
pruning conventions, and the split between stochastic and deterministic-
completion coordinates.  The structural-dynamics chapter already showed that a
filter may approximate the wrong law if it ignores that split.

In the DPF context, the stresses are compounded:
\begin{enumerate}[label=(\roman*)]
  \item the particle-flow map may already be built under closure or local
    linearization assumptions;
  \item the proposal-correction or resampling layer may still introduce finite-
    particle or relaxed-target effects;
  \item invalid or barely valid parameter regions can interact badly with local
    Jacobians, flow stiffness, and repeated HMC evaluations;
  \item a small change in the latent-state law can induce a meaningful change in
    the structural posterior.
\end{enumerate}
Thus nonlinear DSGE models are not a place where one can relax target discipline
simply because the model is already complicated.  They are a place where the
stacking of approximations makes target discipline even more important.

\section{Why MacroFinance models sharpen the issue}
\label{sec:bf-dpf-hmc-macrofinance}

MacroFinance latent-factor systems create a different but equally serious target
challenge.  The main issues include:
\begin{itemize}
  \item higher latent-factor dimension and stronger degeneracy pressure;
  \item mixed-frequency and missing-data structures that complicate repeated
    filtering evaluations;
  \item sharp posterior curvature in volatility and measurement-noise
    directions;
  \item the need for compiled, reproducible long-run evaluation if HMC is to be
    practically useful.
\end{itemize}
A small surrogate drift that appears harmless in a short low-dimensional
experiment can matter materially in a large latent-factor model with long
observation spans and repeated compiled evaluations.  This is why the DPF target
question is not only an abstract theoretical nicety for MacroFinance; it is part
of deciding whether the method remains industrially meaningful.

\section{Relation to surrogate-HMC acceleration}
\label{sec:bf-dpf-hmc-hnn-surrogates}

The DPF path should also be compared to the earlier surrogate-HMC and
Hamiltonian-neural-network acceleration ideas developed in the \texttt{dsge\_hmc} project.  Those
methods try to accelerate sampling by replacing or learning parts of the
Hamiltonian geometry while keeping close track of the target correction issue.

The comparison is instructive.  Learned or surrogate HMC acceleration methods do
not solve a filtering problem; they solve a geometric or computational problem
for an already chosen target.  DPF methods, by contrast, may change the target
itself through flow approximations, resampling relaxations, or learned transport
surrogates.  Thus the two lines of work answer different questions:
\begin{itemize}
  \item HNN/surrogate-HMC asks: can we move through a chosen target more
    efficiently?
  \item DPF target analysis asks: what target have we actually chosen when the
    filter is made differentiable?
\end{itemize}
This means DPF methods and surrogate-HMC methods are not simple substitutes.
DPF concerns the likelihood construction; HNN/surrogate-HMC concerns the
geometry of sampling once a target has been fixed.

\section{BayesFilter recommendation}
\label{sec:bf-dpf-hmc-recommendation}

The rebuilt DPF block now supports a disciplined recommendation.

\begin{enumerate}[label=(\roman*)]
  \item The first justified HMC-relevant DPF development rung is proposal-
    corrected PF-PF, not raw flow and not relaxed resampling.
  \item Soft resampling and OT resampling are mathematically coherent only when
    read as relaxed or surrogate targets, not as hidden exact replacements for
    categorical resampling.
  \item Learned or amortized OT is a further surrogate layer and must be treated
    as such.
  \item For nonlinear DSGE and MacroFinance models, the burden of target
    discipline is higher than in toy systems because structural-state semantics,
    invalid regions, and long-run compiled evaluation amplify target mismatch.
  \item Surrogate-HMC acceleration ideas from the \texttt{dsge\_hmc} project remain relevant as a
    complementary line of work, because they address how to sample a target more
    efficiently once that target has been fixed.
\end{enumerate}

\section{What this chapter does and does not claim}
\label{sec:bf-dpf-hmc-boundary}

This chapter analyzes the DPF ladder as a hierarchy of target constructions for
HMC.  It does \emph{not} yet claim:
\begin{itemize}
  \item that the PF-PF rung has already passed all required variance,
    numerical-stability, or compiled-path gates in BayesFilter;
  \item that a relaxed or learned resampling target is unusable in every
    context;
  \item that nonlinear DSGE or MacroFinance DPF-HMC has already been validated
    in this repository.
\end{itemize}
Its role is narrower and essential: it makes the target question explicit enough
that future experiments and implementations can be judged against a
mathematically clear standard.
